\chapter{Introduction}

\section{Ricci flow}

Let $\mathcal M$ be a smooth closed manifold with Riemannian metric $g$. The
Ricci flow is the evolution equation

\begin{definition}[Ricci flow evolution equation]
\label{topping-intro-ricci-flow-evolution}
\source{topping:page-0007}
\uses{topping-intro-ricci-curvature}
\dcref{ch1:1.1}
\group{ricci-flow}
\lean{Topping.IsRicciFlowOn}
\leanok

\[
\frac{\partial g}{\partial t}=-2\,\Ric(g).
\]
\end{definition}

\begin{remark}
\label{topping-intro-ricci-curvature}
\source{topping:page-0007}
\group{ricci-flow}
\lean{Topping.ricciTensorAt}
\leanok
$\Ric(g)$ denotes the Ricci curvature tensor of $g$.
\end{remark}

In dimension two, a suitable renormalisation of the flow produces a conformal
metric of constant curvature. In general, the flow may develop a finite-time
singularity; surgery removes singular regions before continuation, and the
limiting flow together with the removed pieces can encode the topology.

In harmonic coordinates the Ricci tensor has principal part
\[
R_{ij}=-\frac12\Delta g_{ij}+\text{lower order terms},
\]
whereas at a point $p$ in normal coordinates,
\[
R_{ij}=-\frac32\Delta g_{ij}.
\]
Here $\Delta g_{ij}$ denotes the coordinate expression in these formulas, not
the rough Laplacian of the metric tensor; the latter vanishes because the
metric is parallel.

\section{Special solutions and rescaling}

\subsection{Einstein metrics}

\begin{remark}
\label{topping-intro-einstein-flow}
\source{topping:page-0009}
\group{examples}
If $\Ric(g_0)=\lambda g_0$ for a constant $\lambda\in\mathbb R$, then
\[
g(t)=(1-2\lambda t)g_0
\]
solves Ricci flow with $g(0)=g_0$ wherever the scaling factor is positive.
\end{remark}

\begin{proof}
\sketch Differentiation gives $\partial_tg(t)=-2\lambda g_0$. Uniformly
scaling a metric leaves its Ricci tensor unchanged, so
\[
-2\Ric(g(t))=-2\Ric((1-2\lambda t)g_0)=-2\Ric(g_0)=-2\lambda g_0,
\]
which is the same derivative.
\end{proof}

For the round unit sphere,
\[
\Ric(g_0)=(n-1)g_0,
\qquad
g(t)=(1-2(n-1)t)g_0,
\]
and extinction occurs at
\[
T=\frac{1}{2(n-1)}.
\]
If $g_0$ has constant sectional curvature $-1$, then
$\Ric(g_0)=-(n-1)g_0$, so the corresponding flow expands homothetically for
all time.

\subsection{Ricci solitons}

Let $X(t)$ be a time-dependent smooth vector field with flow $\psi_t$. For a
smooth function $f:\mathcal M\to\mathbb R$, the flow is characterised by:

\begin{definition}[Flow of a time-dependent vector field]
\label{topping-intro-flow-field}
\source{topping:page-0009}
\dcref{ch1:2.1}
\group{ricci-solitons}
\[
X(\psi_t(y),t)f=\frac{\partial}{\partial t}(f\circ\psi_t)(y).
\]
\end{definition}

\begin{proposition}
\label{topping-intro-pushforward-evolution}
\source{topping:page-0010}
\uses{topping-intro-flow-field}
\dcref{ch1:2.1--ch1:2.3}
\group{ricci-solitons}

For a smooth function $\sigma(t)$ and
\[
\hat g(t)=\sigma(t)\psi_t^*(g(t)),
\]
one has
\[
\frac{\partial \hat g}{\partial t}
=\sigma'(t)\psi_t^*(g)
+\sigma(t)\psi_t^*\!\left(\frac{\partial g}{\partial t}\right)
+\sigma(t)\psi_t^*(\mathcal L_Xg).
\]
\end{proposition}

\begin{proof}
\sketch Apply the product rule and the pullback derivative identity
$\frac{d}{dt}\psi_t^*g=\psi_t^*(\mathcal L_Xg)$, with the time dependence of
$g$ contributing the middle term.
\end{proof}

\begin{definition}
\label{topping-intro-ricci-soliton}
\source{topping:page-0010}
\dcref{ch1:2.2--ch1:2.4}
\group{ricci-solitons}
A metric $g_0$, vector field $Y$, and constant $\lambda\in\mathbb R$ define
a Ricci soliton when
\[
-2\Ric(g_0)=\mathcal L_Yg_0-2\lambda g_0.
\]
It is steady, expanding, or shrinking according as $\lambda=0$, $\lambda<0$,
or $\lambda>0$, respectively.
\end{definition}

\begin{definition}
\label{topping-intro-gradient-soliton}
\source{topping:page-0011}
\dcref{ch1:2.3--ch1:2.5}
\group{ricci-solitons}
A Ricci soliton is a gradient Ricci soliton if $Y=\nabla f$ for some smooth
function $f:\mathcal M\to\mathbb R$. Its equation is
\[
\Hess_{g_0}(f)+\Ric(g_0)=\lambda g_0.
\]
\end{definition}

\subsection{Parabolic rescaling}

\begin{remark}
\label{topping-intro-parabolic-rescaling}
\source{topping:page-0012}
\dcref{ch1:2.7--ch1:2.9}
\group{rescaling}
If $g(t)$ is a Ricci flow on $[0,T]$ and $\lambda>0$, then
\[
\hat g(x,t)=\lambda g(x,t/\lambda)
\]
is a Ricci flow on $[0,\lambda T]$. Under this transformation the Ricci tensor
and connection are unchanged, while sectional and scalar curvatures are
multiplied by $\lambda^{-1}$.
\end{remark}

This rescaling permits curvature blow-up analysis by magnifying regions where
curvature becomes large and passing to a pointed limit.

\section{Qualitative behaviour}

In dimension two, $\Ric(g)=Kg$, so negative Gauss curvature expands under the
flow and positive Gauss curvature contracts. In dimension three, a neck pinch
contracts a positively curved spherical cross-section while stretching along
the neck; its blow-up is expected to resemble $S^2\times\mathbb R$, with a
Bryant-soliton model possible for a degenerate asymmetric pinch.

The flow need not converge as $t\to\infty$: hyperbolic metrics expand
homothetically, and long-time analysis can produce thick-thin decompositions
and geometric pieces.

\section{Low-dimensional topology and geometry}

Only closed, oriented manifolds are considered here. Up to diffeomorphism, the
only one-dimensional example is $S^1$. Closed surfaces are classified by genus;
each has a conformally equivalent constant-curvature metric, with curvature
normalised to $1$, $0$, or $-1$. The corresponding universal cover is
$S^2$, $\mathbb R^2$, or $\mathbb H^2$, respectively. Moreover,
\[
\pi_1(S^2)=1,
\qquad
\pi_1(T^2)=\mathbb Z\oplus\mathbb Z,
\]
while the fundamental group of a higher-genus surface is infinite and contains
no subgroup isomorphic to $\mathbb Z\oplus\mathbb Z$.

For closed oriented three-manifolds, prime decomposition reduces the study to
irreducible pieces. The eight model geometries are
\[
S^3,\ \mathbb R^3,\ \mathbb H^3,\ S^2\times\mathbb R,
\ \mathbb H^2\times\mathbb R,\ Nil,\ Sol,
\ \overline{\mathrm{SL}_2}(\mathbb R).
\]
For an irreducible manifold, the expected alternatives include a spherical
space form, a fundamental group containing $\mathbb Z\oplus\mathbb Z$, or a
hyperbolic quotient.

\begin{definition}
\label{topping-intro-geometry}
\source{topping:page-0021}
\dcref{ch1:4.1}
\group{three-manifolds}
A geometry is a simply connected, homogeneous, unimodular Riemannian
manifold.
\end{definition}

Here homogeneous means that any two points are related by an isometry, and
unimodular means that a discrete group of isometries acts with compact
quotient.

\begin{definition}
\label{topping-intro-geometric-manifold}
\source{topping:page-0021}
\dcref{ch1:4.2}
\group{three-manifolds}
A compact manifold $\mathcal M$ is geometric if
$\Int(\mathcal M)=X/\Gamma$ has finite volume, where $X$ is a geometry and
$\Gamma$ is a discrete group of isometries acting freely.
\end{definition}

If a geometric manifold has boundary, its boundary is a union of
incompressible tori.

\begin{conjecture}
\label{topping-intro-thurston-geometrization}
\source{topping:page-0021}
\dcref{ch1:4.3}
\group{three-manifolds}
Every smooth, closed, oriented prime three-manifold is obtained by gluing
finitely many geometric pieces along their boundary tori.
\end{conjecture}

In dimension two, renormalised Ricci flow gives uniformisation. In dimension
three, the geometrisation strategy starts the flow from an arbitrary metric,
uses surgery at neck singularities, and continues the flow; in favourable
positive-Ricci cases extinction yields a spherical space form. In particular,
Hamilton's theorem gives a metric of constant positive sectional curvature on
every closed three-manifold with positive Ricci curvature. General long-time
behaviour is expected to separate into thick and thin regions.
